% ============================================================
\section{Practical Overhead}
\label{sec:overhead}
% ============================================================

Table~\ref{tab:overhead} reports mean encoding time and memory across datasets.
Encoding overhead is not the limiting factor for any QIE.
Amplitude encoding at $35\,\mathrm{ms}$ matches the cost of standardising raw
features and is ${\approx}312\times$ faster than polynomial degree-2
($10.98\,\mathrm{s}$).
Basis encoding at $1.75\,\mathrm{s}$, the most expensive QIE, is still
$6\times$ cheaper than the polynomial expansion.
This is a qualified positive finding: encoding cost is low enough that if the
geometry were right, QIE would be the preferred option on latency-constrained
pipelines.
The bottleneck is the representational mismatch identified in
Section~\ref{sec:analysis}, not speed.

\begin{table}[tbp]
\centering
\caption{Encoding overhead (mean across 10 datasets).}
\label{tab:overhead}
\small
\begin{tabular}{lrrrr}
\toprule
Method & Enc.\ time (ms) & Train time (s) & Enc.\ mem.\ (MB) & Train mem.\ (MB) \\
\midrule
Amplitude (QIE) &     35.2 &    13.1 & 135.7 &   4.8 \\
Angle (QIE)     &    811.3 &   114.2 & 291.0 &   5.6 \\
Basis (QIE)     &  1{,}748 &    39.6 & 436.5 & 345.7 \\
\midrule
Raw linear      &     35.2 &    74.7 &  47.9 &  38.0 \\
RFF             &    193.0 &     2.1 &  65.2 &  24.5 \\
PCA             &  1{,}045 &    14.5 &  29.1 &  15.2 \\
Poly-2          & 10{,}981 &   144.9 & 818.5 & 178.4 \\
\bottomrule
\end{tabular}
\end{table}
